\begin{frame}{例子：一个好题目和一个坏题目的差别}
  \Core{项目工作坊最先要解决的不是“做什么酷”，而是“什么问题有清晰对象、控制点和验证路径”。}
  \begin{block}{坏题目}
    “优化大模型推理性能”过大，既没有对象边界，也没有明确失败条件。
  \end{block}
  \pause
  \begin{block}{好题目}
    “在长短请求混部场景下，验证某种 admission 策略能否改善短请求 TTFT 且不损伤 goodput”。
  \end{block}
  \pause
  \begin{block}{课程结论}
    题目一旦能被写成“对象 + 控制点 + 指标 + 失败条件”，项目才真正开始成立。
  \end{block}
\end{frame}

\begin{frame}{例子：为什么中期汇报最该暴露风险而不是展示进度}
  \Core{中期汇报的价值在于及时暴露错误方向，而不是把已经投入的工程量包装成确定性成果。}
  \begin{block}{好汇报}
    说明当前假设最危险的地方、最可能失败的实验和最需要外界判断的路径。
  \end{block}
  \pause
  \begin{block}{坏汇报}
    只列做了多少代码、跑了多少图，却回避核心假设是否站得住。
  \end{block}
  \pause
  \begin{block}{教学意义}
    中期不是“半成品展示会”，而是“研究方向纠偏会”。
  \end{block}
\end{frame}

\begin{frame}{代码例子：项目提案的一页式模板}
  \Core{工作坊里最好强制学生先写一页提案，再决定是否进入较重实现。}
  \begin{block}{模板片段}
    {\ttfamily\small
    \detokenize{Problem object: ...}\par
    \detokenize{Changed control point: ...}\par
    \detokenize{Expected impact path: ...}\par
    \detokenize{Falsifying result: ...}\par
    }
  \end{block}
  \pause
  \begin{block}{课堂应追问}
    如果其中一项写不出来，说明项目还缺什么？
  \end{block}
  \pause
  \begin{block}{课程结论}
    一页提案的作用不是形式化，而是逼迫项目从“想法”进入“可检验命题”。
  \end{block}
\end{frame}

\begin{frame}{代码例子：最小实验矩阵如何设计}
  \Core{课程项目不需要一开始就大规模扫参，更需要一个能快速否定错误方向的实验矩阵。}
  \begin{block}{实验矩阵草图}
    {\ttfamily\small
    \detokenize{baseline vs method}\par
    \detokenize{short-heavy / mixed / long-heavy workload}\par
    \detokenize{ttft / p99 / goodput / probe metrics}\par
    }
  \end{block}
  \pause
  \begin{block}{课堂应追问}
    为什么先做三种 workload 的小矩阵，往往比直接追求“大而全”更有效？
  \end{block}
  \pause
  \begin{block}{结论}
    研究初期最宝贵的不是更多图，而是更快得到可否定反馈。
  \end{block}
\end{frame}

\begin{frame}{例子：项目失控往往从哪一步开始}
  \Core{大多数课程项目不是在最后失败，而是在最初没有把边界收紧时就已经埋下失败。}
  \begin{block}{失控点一}
    先承诺一个“大一统系统”，没有最小场景和最小机制。
  \end{block}
  \pause
  \begin{block}{失控点二}
    先写很多代码，后补问题定义和实验设计。
  \end{block}
  \pause
  \begin{block}{失控点三}
    只追求正结果，没有提前设计失败也有价值的解释路径。
  \end{block}
\end{frame}

\begin{frame}{课堂练习：把你的题目压缩到一句话}
  \Core{一个能做完的课程项目，往往都能被压缩成一句短而硬的问题句。}
  \begin{block}{练习一}
    用一句话写出你的项目对象、控制点和主指标，不超过 35 个字。
  \end{block}
  \pause
  \begin{block}{练习二}
    再用一句话写出最危险的失败条件。
  \end{block}
  \pause
  \begin{block}{练习三}
    用一句话回答：如果只能保留一个实验，你会做哪个，为什么？
  \end{block}
\end{frame}
